% -*- coding: utf-8 -*-
% !TEX program = xelatex
\documentclass[plain,noanswer,medmath]{../../template/jnuexam}
\usepackage{../../template/kysx}

\begin{document}


\examtitle{2003}{3}


\makepart{填空题}{本题共6小题，每小题4分，满分24分}


\begin{problem}
设$f(x) = \begin{cases}
  x^{\lambda}\cos \frac{1}{x}, & \text{若$x \ne 0$}, \\
                            0, & \text{若$x = 0$},
\end{cases}$其导函数在$x = 0$处连续，则$\lambda$的取值范围是 \fillin{}．
\end{problem}


\begin{problem}
已知曲线$y = x^3 - 3a^2 x + b$与$x$轴相切，则$b^2$可以通过$a$表示为$b^2 =$ \fillin{}．
\end{problem}


\begin{problem}
设$a > 0$，$f(x) = g(x) = \begin{cases}
  a, & \text{若$0 \le x \le 1$}, \\
  0, & \text{其他}.
\end{cases}$而$D$表示全平面，则\goodbreak$I = \iint_D f(x)g(y - x)\dx\dy =$ \fillin{}．
\end{problem}


\begin{problem}
设$n$维向量$\alpha = (a,0, \cdots ,0,a)^T,a < 0$；$E$为$n$阶单位矩阵，
矩阵$A = E - \alpha \alpha^T$，$B = E + \frac{1}{a}\alpha \alpha^T$，其中$A$的逆矩阵为$B$，则$a =$ \fillin{}．
\end{problem}


\begin{problem}
设随机变量$X$和$Y$的相关系数为$0.9$，若$Z = X - 0.4$，则$Y$与$Z$的相关系数为 \fillin{}．
\end{problem}


\begin{problem}
设总体$X$服从参数为$2$的指数分布，$X_1,X_2, \cdots ,X_n$为来自总体$X$的简单随机样本，
则当$n \to \infty$时，$Y_n = \frac{1}{n}\sum_{i = 1}^n X_i^2$依概率收敛于 \fillin{}．
\end{problem}


\makepart{选择题}{本题共6小题，每小题4分，满分24分}


\begin{problem}
设$f(x)$为不恒等于零的奇函数，且$f'(0)$存在，则函数$g(x) = \frac{f(x)}{x}$\pickout{}
\begin{abcd}
\item 在$x = 0$处左极限不存在．
\item 有跳跃间断点$x = 0$．
\item 在$x = 0$处右极限不存在．
\item 有可去间断点$x = 0$．
\end{abcd}
\end{problem}


\begin{problem}
设可微函数$f(x,y)$在点$(x_0,y_0)$取得极小值，则下列结论正确的是\pickout{}
\begin{abcd}
\item $f(x_0,y)$在$y = y_0$处的导数等于零．
\item $f(x_0,y)$在$y = y_0$处的导数大于零．
\item $f(x_0,y)$在$y = y_0$处的导数小于零．
\item $f(x_0,y)$在$y = y_0$处的导数不存在．
\end{abcd}
\end{problem}


\begin{problem}
设$p_n = \frac{a_n + \left| a_n \right|}{2}$，$q_n = \frac{a_n - \left| a_n \right|}{2}$，$n = 1,2, \cdots$，
则下列命题正确的是\pickout{}
\begin{abcd}
\item 若$\sum_{n = 1}^{\infty}a_n$条件收敛，则$\sum_{n = 1}^{\infty}p_n$与$\sum_{n = 1}^{\infty}q_n$都收敛．
\item 若$\sum_{n = 1}^{\infty}a_n$绝对收敛，则$\sum_{n = 1}^{\infty}p_n$与$\sum_{n = 1}^{\infty}q_n$都收敛．
\item 若$\sum_{n = 1}^{\infty}a_n$条件收敛，则$\sum_{n = 1}^{\infty}p_n$与$\sum_{n = 1}^{\infty}q_n$敛散性都不定．
\item 若$\sum_{n = 1}^{\infty}a_n$绝对收敛，则$\sum_{n = 1}^{\infty}p_n$与$\sum_{n = 1}^{\infty}q_n$敛散性都不定．
\end{abcd}
\end{problem}


\begin{problem}
设三阶矩阵$A = \begin{pmatrix}
  a & b & b \\
  b & a & b \\
  b & b & a
\end{pmatrix}$，若$A$的伴随矩阵的秩为$1$，则必有\pickout{}
\begin{abcd}
\item $a = b$或$a + 2b = 0$．
\item $a = b$或$a + 2b \ne 0$．
\item $a \ne b$且$a + 2b = 0$．
\item $a \ne b$且$a + 2b \ne 0$．
\end{abcd}
\end{problem}


\begin{problem}
设$\alpha_1,\alpha_2, \cdots ,\alpha_s$均为$n$维向量，下列结论不正确的是\pickout{}
\begin{abcd}
\item 若对于任意一组不全为零的数$k_1,k_2, \cdots ,k_s$，
      都有$k_1\alpha_1 + k_2\alpha_2 + \cdots + k_s \alpha_s \ne 0$，
      则$\alpha_1,\alpha_2, \cdots ,\alpha_s$线性无关．
\item 若$\alpha_1,\alpha_2, \cdots ,\alpha_s$线性相关，
      则对于任意一组不全为零的数$k_1,k_2, \cdots ,k_s$，
      都有$k_1\alpha_1 + k_2\alpha_2 + \cdots + k_s \alpha_s = 0$．
\item $\alpha_1,\alpha_2, \cdots ,\alpha_s$线性无关的充分必要条件是此向量组的秩为$s$．
\item $\alpha_1,\alpha_2, \cdots ,\alpha_s$线性无关的必要条件是其中任意两个向量线性无关．
\end{abcd}
\end{problem}


\begin{problem}
将一枚硬币独立地掷两次，引进事件：$A_1=$\{掷第一次出现正面\}，$A_2=$\{掷第二次出现正面\}，
$A_3=$\{正、反面各出现一次\}，$A_4=$\{正面出现两次\}，则事件\pickout{}
\begin{abcd}
\item $A_1,A_2,A_3$相互独立．
\item $A_2,A_3,A_4$相互独立．
\item $A_1,A_2,A_3$两两独立．
\item $A_2,A_3,A_4$两两独立．
\end{abcd}
\end{problem}


\makepart{}{本题满分8分}

\begin{problem}
设$f(x) = \frac{1}{\pi x} + \frac{1}{\sin \pi x} - \frac{1}{\pi(1 - x)},x \in\left[\frac{1}{2},1\right)$，
试补充定义$f(1)$使得$f(x)$在$\left[\frac{1}{2},1\right]$上连续．
\end{problem}


\makepart{}{本题满分8分}

\begin{problem}
设$f(u,v)$具有二阶连续偏导数，且满足$\frac{\pd^2 f}{\pd u^2} + \frac{\pd^2 f}{\pd v^2} = 1$，又\goodbreak
$g(x,y) = f[xy,\frac{1}{2}(x^2 - y^2)]$，求$\frac{\pd^2 g}{\pd x^2} + \frac{\pd^2 g}{\pd y^2}$．
\end{problem}


\makepart{}{本题满分8分}

\begin{problem}
计算二重积分
$$I = \iint_D \e^{- (x^2 + y^2 - \pi)}\sin(x^2 + y^2)\dx\dy,$$
其中积分区域$D = \{(x,y)\left| x^2 + y^2 \le \pi \right.\}$．
\end{problem}


\makepart{}{本题满分9分}

\begin{problem}
求幂级数$1 + \sum_{n = 1}^{\infty}(- 1)^n\frac{x^{2n}}{2n}$（$|x| < 1$）的和函数$f(x)$及其极值．
\end{problem}


\makepart{}{本题满分9分}

\begin{problem}
设$F(x) = f(x)g(x)$，其中函数$f(x)$，$g(x)$在$(-\infty , +\infty)$内满足以下条件：\goodbreak
$f'(x) = g(x)$，$g'(x) = f(x)$，且$f(0) = 0$，$f(x) + g(x) = 2\e^x$．
\begin{enumerate}
\item 求$F(x)$所满足的一阶微分方程；
\item 求出$F(x)$的表达式．
\end{enumerate}
\end{problem}


\makepart{}{本题满分8分}

\begin{problem}
设函数$f(x)$在$[0,3]$上连续，在$(0,3)$内可导，且$f(0) + f(1) + f(2) = 3$，$f(3) = 1$．
试证：必存在$\xi \in(0,3)$，使$f'(\xi) = 0$．
\end{problem}


\makepart{}{本题满分13分}

\begin{problem}
已知齐次线性方程组
$$\begin{cases}
  (a_1 + b) x_1 + a_2x_2 + a_3x_3 + \cdots + a_n x_n = 0, \\
  a_1x_1 + (a_2 + b) x_2 + a_3x_3 + \cdots + a_n x_n = 0, \\
  a_1x_1 + a_2x_2 + (a_3 + b) x_3 + \cdots + a_n x_n = 0, \\
  \cdots \cdots \cdots \cdots \cdots \cdots \cdots \cdots \cdots \cdots \\
  a_1x_1 + a_2x_2 + a_3x_3 + \cdots + (a_n + b) x_n = 0,
\end{cases}$$
其中$\sum_{i = 1}^n a_i \ne 0$．试讨论$a_1,a_2, \cdots ,a_n$和$b$满足何种关系时，
\begin{enumerate}
\item 方程组仅有零解；
\item 方程组有非零解．在有非零解时，求此方程组的一个基础解系．
\end{enumerate}
\end{problem}


\makepart{}{本题满分13分}

\begin{problem}
设二次型$f(x_1,x_2,x_3) = X^T AX = ax_1^2 + 2x_2^2 - 2x_3^2 + 2bx_1x_3\;(b > 0)$
中二次型的矩阵$A$的特征值之和为$1$，特征值之积为$-12$．
\begin{enumerate}
\item 求$a$，$b$的值；
\item 利用正交变换将二次型$f$化为标准形，并写出所用的正交变换和对应的正交矩阵．
\end{enumerate}
\end{problem}


\makepart{}{本题满分13分}

\begin{problem}
设随机变量$X$的概率密度为$$f(x) = \begin{cases}
  \frac{1}{3\sqrt[3]{x^2}}, & \text{若$x \in[1,8]$}, \\
                         0, & \text{其他};
\end{cases}$$
$F(x)$是$X$的分布函数．求随机变量$Y = F(X)$的分布函数．
\end{problem}


\makepart{}{本题满分13分}

\begin{problem}
设随机变量$X$与$Y$独立，其中$X$的概率分布为$X\sim \begin{pmatrix}
    1 & 2 \\
  0.3 & 0.7
\end{pmatrix}$，而$Y$的概率密度为$f(y)$，求随机变量$U = X + Y$的概率密度$g(u)$．
\end{problem}


\end{document}
